\section{Related Work}
\label{sec:related}

\paragraph{Rule-guided recurring vulnerability detection.} Closest to us
is the work of Cao et al.~\cite{cao2025vulpa}, VulPA, which also derives
a checkable rule from a known bug and searches a codebase for other
violations. Their rule language, VPDL, compiles to an automaton and
can bind several correlated variables at once, for example checking that
the same object reached both a permission check and a sensitive call; an
IFDS-based interprocedural analysis then walks the whole program. VulPA
targets Java and reports 90 confirmed new findings across 26 applications
and 34 source CVEs. Our rules are intraprocedural, with no alias analysis,
and evaluated on far smaller scale. Protocol~A (Section~\ref{sec:rq5}) is
direct evidence of this gap: its whole-tree run is 0 for 11 precisely
because it lacks the multi-object, ownership-aware correlation VulPA has.

\paragraph{Clone-based recurring vulnerability detection.} Cao et
al.~\cite{cao2025issta} study how well existing recurring-vulnerability
detectors (VUDDY, MVP, MOVERY, V1scan, FIRE) find code \emph{reused} from
a known-vulnerable function. They show these detectors do much worse on
syntactically or semantically modified clones than on exact or renamed
ones. Their notion of \emph{recurring} is a vulnerable function's code
reappearing, by reuse, elsewhere; ours is a different function
independently repeating the same mistake, with no clone relationship (e.g.,
the four KDF functions in Observation~\ref{obs:rq1-conformance}). Tan et al.~\cite{tan2025icse}
address a failure mode specific to clone-based detectors: code that is
similar to a vulnerable function but has itself already been patched, and
which clone-based tools then misflag. \tool checks control flow, not similarity, so it avoids this false-positive
source (though not those in Section~\ref{sec:threats}).

\paragraph{Locating when a vulnerability was introduced.} Guo and
He~\cite{guo2026ndssvic} identify the commit that first introduced a
known vulnerability, by extracting the vulnerable/patched diff's
statement-level patching pattern and matching it backward through commit
history. That asks which commit introduced a known bug, not which other functions
share its pattern now; we take each fix commit as given.

\paragraph{Patch porting.} Wu et al.~\cite{wu2025mystique} port an
existing, known patch to other versions or branches of the same project,
combining semantic and syntactic matching with an LLM. This assumes the target location is known; we find locations not yet known
to need the fix, so the two could be complementary (not evaluated).
